\section{Response degeneracies and the stochastic spectral certificate}\label{sec:supp-r69}
The exact spectral criterion in the main article deliberately assumes a nonreal mode and nonzero contrast. The following closes its real and zero-mode boundaries, without applying a three-mode rank argument where the modes coalesce.

\begin{proposition}[Complete scalar boundary cases]\label{prop:r69-degenerate}
Let $b_n=\tfrac12(1+a\lambda^n)$, with $|a|\le1$, $-1\le\lambda\le1$, and $\lambda^0=1$. Its least stochastic live memory with a prescribed initial label is one if $a=0$ or $\lambda=1$, and otherwise two. In particular this covers the original sinusoidal family at phases $0$ and $\pi$ with $a=\gamma\cos\varphi$, and at $\alpha=0$ with $\lambda=0$.
\end{proposition}
\begin{proof}
In the constant cases a single label with the constant output probability suffices. In every other case $b_0\ne b_1$, whereas a one-label stochastic machine has a constant output probability. For an upper construction use
\[
 Q=\frac12\begin{pmatrix}1+\lambda&1-\lambda\\1-\lambda&1+\lambda\end{pmatrix},
 \qquad d=\frac12\begin{pmatrix}1+a\\1-a\end{pmatrix}.
\]
The constant vector is fixed, and $(1,-1)^{\mathsf T}$ is an eigenvector at $\lambda$. Starting in the first label gives the required response. This includes negative real eigenvalues and both unit-circle endpoints. At $\alpha=0$ the formula gives its one exceptional initial output followed by a constant half output.
\end{proof}

\subsection{Polynomial verification without an initial randomization}
For $z=x+iy$, $y\ne0$, the real variables of \eqref{eq:r69-bilinear} encode a complex eigenvector $v=a+ic$. The equations $Qa=xa-yc$, $Qc=ya+xc$ imply $Qv=zv$, while $v_1=1$ prevents a zero witness. The inequalities $a_j^2+c_j^2\le1$ ensure legal decoder probabilities for every $0<\gamma\le1$ and every fixed phase. Conversely a maximal-modulus eigenvector coordinate can always be normalized this way. Thus the certificate does not assume a special circulant form, a full-support stationary distribution, a reversible chain or a hidden preparation mixture.

With algebraic $x,y,\gamma,\cos\varphi,\sin\varphi$, every nonempty fixed-size certificate has an algebraic sample by real-closed-field elimination. A dense transition table has $M^2$ entries and a binary decoder has $M$ probabilities, in addition to the retained $\lceil\log_2M\rceil$ state bits. Algebraic-number representations and exact sampling have their own costs. The cubic-root theorem does not silently grant a polynomial bit implementation of an arbitrary real Diophantine rotation.

\subsection{Constructing a directional circulant}
Let $q\ge3$ and place $\vartheta$ between adjacent root angles $2\pi j/q$ and $2\pi(j+1)/q$, choosing a representative on the circle. Write $\beta=2\pi/q$ and $\psi=\vartheta-2\pi j/q\in[0,\beta]$. If $\alpha\le R_q(\vartheta)$, set
\[
 a=\alpha\frac{\sin(\beta-\psi)}{\sin\beta},\qquad
 b=\alpha\frac{\sin\psi}{\sin\beta},\qquad c=1-a-b.
\]
All three are nonnegative. The matrix $Q=aC_q^j+bC_q^{j+1}+cJ$ is stochastic and has eigenvector $v_l=e^{2\pi il/q}$ at $z=\alpha e^{i\vartheta}$. Its decoder is $d_l=(1+\gamma\cos(\varphi+2\pi l/q))/2$. This is a directly specified program at a fixed phase, not an abstract existence proof requiring the physical state to be observed. A root-angle endpoint is permitted and may reduce the actual minimal size. The construction is sufficient; it is not asserted to describe every part of the Karpelevi\v c boundary.

The cyclic eigenvector gives the all-time certificate $Qv=zv$. Checking a finite list of output probabilities without that identity would not certify exact realization. The new diagnostic program tests both the intertwining identity and independent raw qubit evolution, and labels the latter tests as finite probes.

\subsection{Constants and quantifiers in the arithmetic theorem}
For $h\le\pi/3$ and $0\le\delta\le h$, the inequalities used in the lower proof are
\[
 \sin\delta\ge\frac34\delta,\qquad
 1-\cos\delta\le\frac12\delta^2,\qquad
 \tan h\ge h,\qquad
 \cos\delta+\sin\delta\tan h\le\sec h\le2.
\]
The first follows, for example, from concavity of sine and
$\sin(\pi/3)/(\pi/3)=3\sqrt3/(2\pi)>3/4$. Thus there is no numerical fitting in the constant $16$ of the cubic-root lower. The exact first-sector law proves the worst-phase lower, not a relaxation to a larger convex hull of all low-order roots. In contrast, the fixed-phase upper uses a polygon whose size is a Diophantine denominator; the two constructions have different quantifiers.

For the golden quadratic phase $x=(3-\sqrt5)/2$, the integer-polynomial proof in the main text gives a global constant $c=1/3$, valid at every positive denominator. Inspecting a finite continued-fraction prefix would not give that global conclusion. For $(3+4i)/5$ the Gaussian-integer argument gives the displayed logarithmic lower but not a badly-approximable certificate. These facts prevent transferring the cubic-root exponent to a different phase by numerical resemblance.

The exact-memory count is discontinuous at vanishing contrast: all nonzero contrasts have the same spectral obstruction, but zero contrast is constant. This statement concerns equality of every response coefficient. Approximate risk at a fixed tolerance can of course become insensitive to small contrast; that is a different optimization problem and is not contradicted by the exact criterion.

\subsection{Relation to classical results and to other resource curves}
The stochastic sector inequality, Karpelevi\v c regions and first-sector lazy-cycle classification are classical; their pertinent statements are credited in the main text and proved there in the form used. The noise-dependent classical-memory bounds in Fanizza--Lumbreras--Winter concern processes and quantifiers distinct from the current one-response squared task. The present argument neither claims priority for the unit-circle obstruction nor uses it as a substitute for the matched arithmetic rate.

The protected system in the squared embedding has dimension four, independent of the observer labels. The instrument-level state-positive simulator of Section~\ref{sec:interior-instruments} is a separate construction: it preserves every declared terminal effect and adaptive experiment, whereas the current stochastic response realization preserves only the fixed optimal task kernel. The old two-label squared-sign frontier and delayed-release timing frontier likewise have different objective functions. Their proofs remain current and no general-$M$ risk formula for either is inferred from the new exact-response law.
